\documentclass[10pt,a4paper]{amsart}
\usepackage[margin=19mm]{geometry}
\usepackage{amsmath,amssymb,mathtools}
\usepackage[scr=rsfs,cal=euler]{mathalfa}
\usepackage{xfrac}
\usepackage[unicode,hidelinks,bookmarks=false]{hyperref}
\newcommand{\ie}{i.e.\ }
\newcommand{\cf}{cf.\ }
\newcommand{\resp}{respectively\ }
\newcommand{\Subsection}{\subsection}
\providecommand{\nobreakdash}{\nobreak\hbox{-}}
\setlength{\emergencystretch}{2em}
\multlinegap0pt
\usepackage{amssymb}
\usepackage{bm}
\usepackage{tikz}
\usepackage{algorithm}\usepackage{algpseudocode}
\newtheorem{theorem}{Theorem}
\title{Quantum Viterbi Algorithm}
\author{Luigi Accardi, Abdessatar Souissi, El Gheteb Soueidi, Farrukh Mukhamedov, Mohamed Rhaima}
\date{Source: arXiv:2605.18912v1 (CC BY 4.0)}
\begin{document}
\maketitle

\noindent\emph{Source and licence.} Quantum Viterbi Algorithm. Version arXiv:2605.18912v1 (CC BY 4.0), distributed by arXiv under the Creative Commons Attribution 4.0 International licence, http://creativecommons.org/licenses/by/4.0/, as recorded in the arXiv licence metadata for this exact version and shown on the abstract page https://arxiv.org/abs/2605.18912v1. Full licence notice: \texttt{licenses/arxiv-2605.18912-CC-BY-4.0.txt}.\par\smallskip

\noindent\textit{Editorial scope.} Counted unit: the article's theorem thm:coherent-optimal-m0 (source0.tex lines 748--757). The article's complete original proof of it is delivered with it (source0.tex lines 759--855, one \texttt{proof} environment of 97 lines). No mathematics is written, repaired or re-derived: the statement and the proof are the article's own lines, in the article's own notation, and no retained line is substituted or re-flowed.  No further source-local context is needed: every cross-reference of every retained line resolves inside the two delivered spans.  Nothing was omitted from any retained span, so the package carries no omission record.  No retained line contains a citation, so the wrapper carries no bibliography and no bibliography entry.  No retained line contains a figure environment, an image-inclusion command or drawing code, so no image is delivered and the package declares no asset dependency.  Editorial formatting only, in addition: an AMS \texttt{amsart} preamble that restates the article's own declarations and macros used by the retained lines, namely the environments \texttt{theorem} on separate counters, together with the standard packages those lines need.  The retained environments print in this document as Theorem 1 in order of appearance, whereas the article prints the same items with the numbering of its own full document.  The complete native source of the article as distributed in the arXiv e-print for this exact version is bundled unchanged as \texttt{sources/200882/source0.tex}.
\begin{theorem}[Single-step quantum advantage]
\label{thm:coherent-optimal-m0}
There exist physically admissible parameters \((\eta_0,p_0,\phi_0)\), an observation \(q_0\), a continuation operator \(V_1\), and an initial functional \(\varphi_{H,0}\) such that the one-step functional \(F_0\) satisfies
\[
\sup_{p_0\in\mathcal{P}_{1}(\mathcal{H})} F_0(p_0)
\;>\;
\sup_{p_0\in D_H} F_0(p_0),
\]
where \(D_H=\{|0\rangle\langle0|,|1\rangle\langle1|\}\) denotes the classical subset of diagonal pure effects. Moreover, any hidden effect \(p_0^*\) achieving the supremum on the left must be coherent, in the sense that it has nonzero off-diagonal entries in the computational basis.
\end{theorem}

\begin{proof}
We argue by explicit construction. Choose parameters as follows: set dephasing \(\theta_0 = 0\) so that the hidden transition at time \(0\) is purely unitary; choose \(\phi_0 = \pi/4\), giving \(U_0 = e^{-i(\pi/4)\sigma_x}\); set the measurement strength \(\eta_0 = 1/2\), so that
\[
L_0^{+} = \sqrt{\tfrac{3}{2}}\,|0\rangle\langle0| + \sqrt{\tfrac{1}{2}}\,|1\rangle\langle1|;
\]
fix the observation to be \(q_0 = |+\rangle\langle+|\), where \(|+\rangle = (|0\rangle+|1\rangle)/\sqrt2\); choose the continuation operator to be \(V_1 =  I_2\); and take the initial functional \(\varphi_{H,0}(X) = \mathrm{Tr}(aX)\) with \(a = |\phi_0\rangle\langle\phi_0|\) and \(|\phi_0\rangle = (|0\rangle+|1\rangle)/\sqrt2\).

With \(\theta_0=0\) and \(\mathrm{Tr}(V_1)=1\) (normalized trace), the hidden transition simplifies to \(\mathcal{E}_{H;0}(X_1\otimes V_1) = U_0^\dagger X_1 U_0\). Using that \(U_0|\phi_0\rangle = e^{-i\pi/4}|\phi_0\rangle\), one finds \(U_0 a U_0^\dagger = a\). Consequently, the one-step functional reduces to
\[
F_0(p_0) = \mathrm{Tr}\bigl(a\,\mathcal{E}_{H,O;0}(p_0\otimes q_0)\bigr),
\]
and, since \(q_0 = |+\rangle\langle+|\) selects the \(L_0^{+}\) outcome, we have \(\mathcal{E}_{H,O;0}(p_0\otimes q_0) = L_0^{+} p_0 L_0^{+}\).

Let \(p_0 = |\psi\rangle\langle\psi|\) with \(|\psi\rangle = \alpha|0\rangle + \beta|1\rangle\) and \(|\alpha|^2+|\beta|^2=1\). A direct matrix calculation gives
\[
L_0^{+} p_0 L_0^{+}
  =
  \begin{pmatrix}
    \tfrac{3}{2}|\alpha|^2 & \tfrac{\sqrt{3}}{2}\,\alpha\overline{\beta} \\
    \tfrac{\sqrt{3}}{2}\,\overline{\alpha}\beta & \tfrac{1}{2}|\beta|^2
  \end{pmatrix}.
\]
With \(a = \tfrac12\begin{pmatrix}1&1\\[2pt]1&1\end{pmatrix}\), this yields
\[
F_0(p_0)
  = \tfrac{3}{8}|\alpha|^2 + \tfrac{1}{8}|\beta|^2
    + \tfrac{\sqrt{3}}{8}(\alpha\overline{\beta} + \overline{\alpha}\beta).
\]
On the classical subset \(D_H\), one finds \(F_0(|0\rangle\langle0|) = 3/8\) and \(F_0(|1\rangle\langle1|) = 1/8\), so
\[
\sup_{p_0\in D_H} F_0(p_0) = 3/8.
\]

Let
\[
x =
\begin{pmatrix}
\alpha \\
\beta
\end{pmatrix}
\in \mathbb{C}^2,
\quad \text{with } \|x\|^2 = |\alpha|^2 + |\beta|^2 = 1.
\]
Then \(F_0(p_0)\) can be written as a quadratic form
\[
F_0(p_0)=G(\alpha,\beta) = x^* A x,
\]
where
\[
A =
\frac{1}{8}
\begin{pmatrix}
3 & \sqrt{3} \\
\sqrt{3} & 1
\end{pmatrix}.
\]
The matrix \(A\) is Hermitian, hence it admits a spectral decomposition with real eigenvalues  \(1/2\) and \(0\). Therefore, for any unit vector \(x \in \mathbb{C}^2\),
\[
x^* A x = \lambda_1 |c_1|^2 + \lambda_2 |c_2|^2
\]
where \(\lambda_1,\lambda_2\) are the eigenvalues of \(A\), and
\[
|c_1|^2 + |c_2|^2 = 1.
\]
Thus \(x^* A x\) is a convex combination of the eigenvalues, and we obtain
\[
0
\le x^* A x
\le \tfrac{1}{2}.
\]
Consequently,
\[
\sup_{\|x\|=1} x^* A x = 1/2.
\]

This yields
\[
\sup_{|\alpha|^2 + |\beta|^2 = 1} G(\alpha,\beta) = \frac{1}{2}.
\]

For the coherent state
$$
|\psi^*\rangle
=
\frac{\sqrt{3}}{2}|0\rangle + \frac{1}{2}|1\rangle,
\qquad
p_0^* = |\psi^*\rangle\langle\psi^*|
$$
we have
\[
\sup_{p_0\in\mathcal{P}_1(\mathcal{H})} F_0(p_0)=F_0(p_0^*)
  = \tfrac{1}{2} > \sup_{p_0\in D_H} F_0(p_0) = 3/8.
\]

as claimed.

\end{proof}

\end{document}
